\section{Tagged PDF 与替代文本}
\label{sec:manual-accessibility}
在支持 tagging 的 LaTeX 发行版中，在 documentclass 之前使用 DocumentMetadata 开启标记。
宏包为独立图形自动生成英文替代文本，包括原始数值、单位、量程、参考线、误差距离、
对比两端与差值、堆叠分段与正负小计，以及趋势中的每个观测和缺失位置。
区间描述保留上下界与显式点估计，不推断置信水平。图例描述类别名称。
表格测量过程不生成语义标签，实际图形单元格只生成一次描述。

\begin{codeonly}
\DocumentMetadata{tagging=on,pdfstandard=ua-2,lang=zh-CN}
\documentclass{ctexart}
\usepackage{gradbars}
\begin{document}
\gradbar{72}

\gradbar[alt={准确率为百分之七十二，满分为百分之百。}]{72}

\graddumbbell[compare label=delta]{58}{76}
\end{document}
\end{codeonly}

\key{alt={描述}} 覆盖自动文本；空值恢复自动生成。
原有 \key{/tikz/alt={描述}} 写法继续可用。

\key{accessibility} 可选 auto、artifact 或 off。
auto 是默认值，为图形生成描述；artifact 将图形作为装饰，
仅用于周围文字已完整表达其信息的情况；off 将标记设置交由外层控制。

未开启 tagging 时这些选项不改变视觉输出。自定义描述应解释单位、类别及图形在当前文档中的含义。

在已有 tikzpicture 内调用宏包时不会创建独立图形标签，应为外层图提供完整替代文本。
自动描述使用原始数据，不随 precision 舍入；text 仅修改视觉标签，语义描述需通过 alt 覆盖。
该支持不等于整份文档已符合 PDF/UA；还需检查表头、阅读顺序、字体及其他内容的标记。
